\chapter{svndumpfilter Reference---Subversion History
Filtering}\label{svn.ref.svndumpfilter}

\texttt{svndumpfilter} is a command-line utility for removing history
from a Subversion dump file by either excluding or including paths
beginning with one or more named prefixes. For details, see
\hyperref[svn.reposadmin.maint.tk.svndumpfilter]{svndumpfilter}.

Options in \texttt{svndumpfilter} are global, just as they are in
\texttt{svn} and \texttt{svnadmin}:

\protect\phantomsection\label{svn.ref.svndumpfilter.sw}{}

\begin{description}
\item[\protect\phantomsection\label{svn.ref.svndumpfilter.sw.drop_empty_revs}{}\texttt{-\/-drop-empty-revs}]
If the current filtering invocation causes any revision to be empty
(i.e., the revision causes no change to the repository), removes these
revisions from the final dump file.
\item[\protect\phantomsection\label{svn.ref.svndumpfilter.sw.drop_all_empty_revs}{}\texttt{-\/-drop-all-empty-revs}]
Removes all empty revisions found in the dumpstream from the final dump
file (except revision 0).
\item[\protect\phantomsection\label{svn.ref.svndumpfilter.sw.pattern}{}\texttt{-\/-pattern}]
Treat the path prefixes provided to the filtering commands as file glob
patterns rather than explicit path substrings.
\item[\protect\phantomsection\label{svn.ref.svndumpfilter.sw.renumber_revs}{}\texttt{-\/-renumber-revs}]
Renumbers revisions that remain after filtering.
\item[\protect\phantomsection\label{svn.ref.svndumpfilter.sw.skip_missing_merge_sources}{}\texttt{-\/-skip-missing-merge-sources}]
Skips merge sources that have been removed as part of the filtering.
Without this option, \texttt{svndumpfilter} will exit with an error if
the merge source for a retained path is removed by filtering.
\item[\protect\phantomsection\label{svn.ref.svndumpfilter.sw.preserve_revprops}{}\texttt{-\/-preserve-revprops}]
If all nodes in a revision are removed by filtering and
\texttt{-\/-drop-empty-revs} is not passed, the default behavior of
\texttt{svndumpfilter} is to remove all revision properties except for
the date and the log message (which will merely indicate that the
revision is empty). Passing this option will preserve existing revision
properties (which may or may not make sense since the related content is
no longer present in the dump file).
\item[\protect\phantomsection\label{svn.ref.svndumpfilter.sw.targets}{}\texttt{-\/-targets}
\textless FILENAME\textgreater{}]
Instructs \texttt{svndumpfilter} to read additional path prefixes---one
per line---from the file located at \textless FILENAME\textgreater. This
is especially useful for complex filtering operations which require more
prefixes than the operating system allows to be specified on a single
command line.
\item[\protect\phantomsection\label{svn.ref.svndumpfilter.sw.quiet}{}\texttt{-\/-quiet}]
Does not display filtering statistics.
\end{description}

\section{svndumpfilter
exclude}\label{svn.ref.svndumpfilter.commands.c.exclude}

\emph{Filter out nodes with given prefixes from the dump stream.}

\textbf{Synopsis}

\texttt{svndumpfilter\ exclude\ PATH\_PREFIX...}

\subsection{Description}

This can be used to exclude nodes that begin with one or more
\textless PATH\_PREFIX\textgreater es from a filtered dump file.

\subsection{Options}

\begin{verbatim}
--drop-empty-revs
--drop-all-empty-revs
--pattern
--preserve-revprops
--quiet
--renumber-revs
--skip-missing-merge-sources
--targets FILENAME
\end{verbatim}

\subsection{Examples}

If we have a dump file from a repository with a number of different
picnic-related directories in it, but we want to keep everything
\emph{except} the \texttt{sandwiches} part of the repository,
we\textquotesingle ll exclude only that path:

\begin{verbatim}
$ svndumpfilter exclude sandwiches < dumpfile > filtered-dumpfile
Excluding prefixes:
   '/sandwiches'

Revision 0 committed as 0.
Revision 1 committed as 1.
Revision 2 committed as 2.
Revision 3 committed as 3.
Revision 4 committed as 4.

Dropped 1 node(s):
   '/sandwiches'
$
\end{verbatim}

Beginning in Subversion 1.7, \texttt{svndumpfilter} can optionally treat
the \textless PATH\_PREFIX\textgreater s not merely as explicit
substrings, but as file patterns instead. So, for example, if you wished
to filter out paths which ended with \texttt{.OLD}, you would do the
following:

\begin{verbatim}
$ svndumpfilter exclude --pattern "*.OLD" < dumpfile > filtered-dumpfile
Excluding prefix patterns:
   '/*.OLD'

Revision 0 committed as 0.
Revision 1 committed as 1.
Revision 2 committed as 2.
Revision 3 committed as 3.
Revision 4 committed as 4.

Dropped 3 node(s):
   '/condiments/salt.OLD'
   '/condiments/pepper.OLD'
   '/toppings/cheese.OLD'
$
\end{verbatim}

\section{svndumpfilter
include}\label{svn.ref.svndumpfilter.commands.c.include}

\emph{Filter out nodes without given prefixes from dump stream.}

\textbf{Synopsis}

\texttt{svndumpfilter\ include\ PATH\_PREFIX...}

\subsection{Description}

Can be used to include nodes that begin with one or more
\textless PATH\_PREFIX\textgreater es in a filtered dump file (thus
excluding all other paths).

\subsection{Options}

\begin{verbatim}
--drop-empty-revs
--drop-all-empty-revs
--pattern
--preserve-revprops
--quiet
--renumber-revs
--skip-missing-merge-sources
--targets FILENAME
\end{verbatim}

\subsection{Example}

If we have a dump file from a repository with a number of different
picnic-related directories in it, but want to keep only the
\texttt{sandwiches} part of the repository, we\textquotesingle ll
include only that path:

\begin{verbatim}
$ svndumpfilter include sandwiches < dumpfile > filtered-dumpfile
Including prefixes:
   '/sandwiches'

Revision 0 committed as 0.
Revision 1 committed as 1.
Revision 2 committed as 2.
Revision 3 committed as 3.
Revision 4 committed as 4.

Dropped 12 node(s):
   '/condiments'
   '/condiments/pepper'
   '/condiments/pepper.OLD'
   '/condiments/salt'
   '/condiments/salt.OLD'
   '/drinks'
   '/snacks'
   '/supplies'
   '/toppings'
   '/toppings/cheese'
   '/toppings/cheese.OLD'
   '/toppings/lettuce'
$
\end{verbatim}

Beginning in Subversion 1.7, \texttt{svndumpfilter} can optionally treat
the \textless PATH\_PREFIX\textgreater s not merely as explicit
substrings, but as file patterns instead. So, for example, if you wished
to include only paths which ended with \texttt{ks}, you would do the
following:

\begin{verbatim}
$ svndumpfilter include --pattern "*ks" < dumpfile > filtered-dumpfile
Including prefix patterns:
   '/*ks'

Revision 0 committed as 0.
Revision 1 committed as 1.
Revision 2 committed as 2.
Revision 3 committed as 3.
Revision 4 committed as 4.

Dropped 11 node(s):
   '/condiments'
   '/condiments/pepper'
   '/condiments/pepper.OLD'
   '/condiments/salt'
   '/condiments/salt.OLD'
   '/sandwiches'
   '/supplies'
   '/toppings'
   '/toppings/cheese'
   '/toppings/cheese.OLD'
   '/toppings/lettuce'
$
\end{verbatim}

\section{svndumpfilter
help}\label{svn.ref.svndumpfilter.commands.c.help}

\emph{Help!}

\textbf{Synopsis}

\texttt{svndumpfilter\ help\ {[}SUBCOMMAND...{]}}

\subsection{Description}

Displays the help message for \texttt{svndumpfilter}. Unlike other help
commands documented in this chapter, there is no witty commentary for
this help command. The authors of this book deeply regret the omission.

\subsection{Options}

None

